% =============================================================================
% Preambles/header.tex — shared LaTeX/Beamer preamble for 03_AIWTPGap
%
% Korean experimental-economics proposal. Compile with XeLaTeX:
%     % =============================================================================
% Preambles/header.tex — shared LaTeX/Beamer preamble for 03_AIWTPGap
%
% Korean experimental-economics proposal. Compile with XeLaTeX:
%     % =============================================================================
% Preambles/header.tex — shared LaTeX/Beamer preamble for 03_AIWTPGap
%
% Korean experimental-economics proposal. Compile with XeLaTeX:
%     % =============================================================================
% Preambles/header.tex — shared LaTeX/Beamer preamble for 03_AIWTPGap
%
% Korean experimental-economics proposal. Compile with XeLaTeX:
%     \input{header}
% with TEXINPUTS=../Preambles:$TEXINPUTS (the /compile-latex skill does this).
%
% Required font: KoPubWorld Dotum (Medium / Bold). Install from
% https://www.kopus.or.kr/biz/electronic/font.do (free for academic use).
% Fallback options documented under "Korean font setup" below.
% =============================================================================

% -----------------------------------------------------------------------------
% Engine detection + encoding
%
% XeLaTeX is required (Korean rendering via xeCJK depends on it). Under pdfTeX
% xeCJK is unavailable, so we fall back to inputenc and emit a clear warning.
% -----------------------------------------------------------------------------
\usepackage{iftex}
\ifPDFTeX
  \usepackage[utf8]{inputenc}
  \PackageWarningNoLine{header}{%
    pdfTeX detected --- Korean text will not render. Compile with XeLaTeX.}
\fi

\usepackage{amsmath, amssymb, amsthm}
\usepackage{booktabs}
\usepackage{graphicx}

% -----------------------------------------------------------------------------
% Korean font setup (XeLaTeX only)
%
% Loads xeCJK and KoPubWorld Dotum. The fonts live under ~/Library/Fonts/ on
% macOS. We pass file names (not family names) because the family name is
% registered as Korean ("KoPubWorld돋움체") and fontspec is more robust with
% explicit file references.
%
% Fallback ladder if KoPubWorld Dotum is unavailable: comment out the block
% below and uncomment ONE of the alternatives.
% -----------------------------------------------------------------------------
\ifXeTeX
  \usepackage{xeCJK}
  \xeCJKsetup{%
    CJKspace=true,           % automatic CJK-Latin spacing
    AutoFakeSlant=0.18,      % italic-style fallback when no italic cut exists
    CJKecglue={\hskip 0.18em plus 0.04em minus 0.04em}%
  }

  % --- KoPubWorld Dotum (primary) -------------------------------------------
  % macOS Font Book registers these by PostScript name (no spaces). XeLaTeX's
  % fontspec finds them via the system font registry without needing a
  % Path= override. If your machine doesn't find them, install KoPubWorld
  % Dotum (Medium / Bold / Light) into ~/Library/Fonts/ and rebuild the
  % font cache, or comment this block out and uncomment the Apple SD Gothic
  % Neo fallback below.
  \setCJKmainfont[%
    BoldFont=KoPubWorldDotumBold,
    AutoFakeSlant=0.18
  ]{KoPubWorldDotumMedium}
  \setCJKsansfont[%
    BoldFont=KoPubWorldDotumBold,
    AutoFakeSlant=0.18
  ]{KoPubWorldDotumMedium}
  \setCJKmonofont{KoPubWorldDotumLight}

  % --- Apple SD Gothic Neo fallback (uncomment if KoPub unavailable) --------
  % \setCJKmainfont{Apple SD Gothic Neo}
  % \setCJKsansfont{Apple SD Gothic Neo}
  % \setCJKmonofont{Apple SD Gothic Neo}

  % --- Pretendard fallback (uncomment if installed) -------------------------
  % \setCJKmainfont{Pretendard}
  % \setCJKsansfont{Pretendard}
  % \setCJKmonofont{Pretendard}

  % Slightly looser line spread improves Hangul readability on slides
  \linespread{1.15}
\fi

% -----------------------------------------------------------------------------
% PALETTE — Beamer theme + TikZ-snippet colors.
% Load xcolor and define colors BEFORE anything that references them
% (notably hyperref, hypersetup, and the Beamer color assignments).
% -----------------------------------------------------------------------------
\usepackage{xcolor}

\definecolor{primary-blue}{HTML}{012169}      % headings, accents
\definecolor{primary-gold}{HTML}{B9975B}      % emphasis, borders
\definecolor{highlight-yellow}{HTML}{F2A900}  % markers, alerts
\definecolor{light-bg}{HTML}{E8EDF5}          % title / section backgrounds
\definecolor{jet}{HTML}{1A1A1A}               % body text

% Semantic colors (used in diagrams and annotations)
\definecolor{positive}{HTML}{15803D}          % good / observed / identified
\definecolor{negative}{HTML}{B91C1C}          % bad / problematic / bias
\definecolor{neutral}{HTML}{525252}           % reference / context

% Highlight accents (match .hi-* classes in the SCSS)
\definecolor{hi-slate}{HTML}{314F4F}
\definecolor{hi-green}{HTML}{2E8B57}
\definecolor{hi-red}{HTML}{C41E3A}

% Semantic aliases so skill templates and snippets can write
% \textcolor{accent}{...} without hard-coding "primary-blue".
\colorlet{accent}{primary-blue}
\colorlet{accent2}{primary-gold}

% -----------------------------------------------------------------------------
% Hyperref — loaded AFTER xcolor + palette so link colors resolve.
% -----------------------------------------------------------------------------
\usepackage{hyperref}
\hypersetup{colorlinks=true, linkcolor=primary-blue, urlcolor=primary-blue,
            citecolor=primary-blue, pdfborder={0 0 0}}

% -----------------------------------------------------------------------------
% Beamer theme assignments (only apply under Beamer).
%
% We detect Beamer via \@ifclassloaded{beamer}, which is the documented,
% non-internal way. This must be wrapped in \makeatletter/\makeatother
% because \@ifclassloaded contains an @-catcode character.
% -----------------------------------------------------------------------------
\makeatletter
\@ifclassloaded{beamer}{%
  \usefonttheme{professionalfonts}
  \setbeamercolor{structure}{fg=primary-blue}
  \setbeamercolor{frametitle}{fg=primary-blue}
  \setbeamercolor{title}{fg=primary-blue}
  \setbeamercolor{subtitle}{fg=primary-gold}
  \setbeamercolor{author}{fg=primary-blue}
  \setbeamercolor{institute}{fg=neutral}
  \setbeamercolor{date}{fg=primary-gold}
  \setbeamercolor{itemize item}{fg=primary-blue}
  \setbeamercolor{itemize subitem}{fg=primary-gold}
  \setbeamercolor{alerted text}{fg=primary-gold}
  \setbeamercolor{block title}{fg=white, bg=primary-blue}
  \setbeamercolor{block body}{bg=light-bg}

  % Minimal footer: page X / N only, no navigation clutter
  \setbeamertemplate{navigation symbols}{}
  \setbeamertemplate{footline}{%
    \hfill{\color{neutral}\footnotesize\insertframenumber/\inserttotalframenumber}\hspace{0.7em}\vspace{0.3em}}
}{}
\makeatother

% -----------------------------------------------------------------------------
% TikZ setup — libraries the snippet gallery uses
% -----------------------------------------------------------------------------
\usepackage{tikz}
\usetikzlibrary{arrows.meta, positioning, calc, decorations.pathreplacing,
                fit, shapes.geometric, backgrounds}

% Shared TikZ styles — snippets and hand-written diagrams can pick these up
% via `\begin{tikzpicture}[every node/.style={...}]` or by naming specific
% styles (e.g., `\node[dag-node] (X) at (0,0) {X};`).
\tikzset{
  % Node styles for causal DAGs and similar
  dag-node/.style = {
    draw=primary-blue, fill=light-bg, rounded corners=2pt,
    minimum width=1.2cm, minimum height=0.9cm, align=center, font=\small
  },
  decision-node/.style = {
    draw=primary-gold, fill=white, diamond, aspect=1.8,
    minimum width=1.8cm, minimum height=1.2cm, align=center, font=\small,
    inner sep=1pt
  },
  % Edge styles — observed vs counterfactual
  observed-edge/.style = {-{Latex[length=2.5mm]}, thick, draw=jet},
  counterfactual-edge/.style = {-{Latex[length=2.5mm]}, thick, draw=neutral, dashed},
  confound-edge/.style = {-{Latex[length=2.5mm]}, thick, draw=negative, dashed},
  % Data-point styles for plots
  observed-dot/.style = {fill=primary-blue, draw=primary-blue, circle, inner sep=1.2pt},
  counterfactual-dot/.style = {fill=white, draw=neutral, circle, inner sep=1.2pt}
}

% -----------------------------------------------------------------------------
% Convenience macros
% -----------------------------------------------------------------------------
\newcommand{\muted}[1]{\textcolor{neutral}{#1}}
\newcommand{\key}[1]{\textcolor{primary-gold}{\textbf{#1}}}
\newcommand{\good}[1]{\textcolor{positive}{#1}}
\newcommand{\bad}[1]{\textcolor{negative}{#1}}

% Standout frame macro: full-bleed dark background for section transitions.
% Expand: \transitionslide{Your transition title}
\newcommand{\transitionslide}[1]{%
  \begingroup
  \setbeamercolor{background canvas}{bg=primary-blue}
  \begin{frame}[plain]
    \centering
    \vfill
    {\usebeamerfont{title}\color{white}\Large #1\par}
    \vfill
  \end{frame}
  \endgroup
}

% with TEXINPUTS=../Preambles:$TEXINPUTS (the /compile-latex skill does this).
%
% Required font: KoPubWorld Dotum (Medium / Bold). Install from
% https://www.kopus.or.kr/biz/electronic/font.do (free for academic use).
% Fallback options documented under "Korean font setup" below.
% =============================================================================

% -----------------------------------------------------------------------------
% Engine detection + encoding
%
% XeLaTeX is required (Korean rendering via xeCJK depends on it). Under pdfTeX
% xeCJK is unavailable, so we fall back to inputenc and emit a clear warning.
% -----------------------------------------------------------------------------
\usepackage{iftex}
\ifPDFTeX
  \usepackage[utf8]{inputenc}
  \PackageWarningNoLine{header}{%
    pdfTeX detected --- Korean text will not render. Compile with XeLaTeX.}
\fi

\usepackage{amsmath, amssymb, amsthm}
\usepackage{booktabs}
\usepackage{graphicx}

% -----------------------------------------------------------------------------
% Korean font setup (XeLaTeX only)
%
% Loads xeCJK and KoPubWorld Dotum. The fonts live under ~/Library/Fonts/ on
% macOS. We pass file names (not family names) because the family name is
% registered as Korean ("KoPubWorld돋움체") and fontspec is more robust with
% explicit file references.
%
% Fallback ladder if KoPubWorld Dotum is unavailable: comment out the block
% below and uncomment ONE of the alternatives.
% -----------------------------------------------------------------------------
\ifXeTeX
  \usepackage{xeCJK}
  \xeCJKsetup{%
    CJKspace=true,           % automatic CJK-Latin spacing
    AutoFakeSlant=0.18,      % italic-style fallback when no italic cut exists
    CJKecglue={\hskip 0.18em plus 0.04em minus 0.04em}%
  }

  % --- KoPubWorld Dotum (primary) -------------------------------------------
  % macOS Font Book registers these by PostScript name (no spaces). XeLaTeX's
  % fontspec finds them via the system font registry without needing a
  % Path= override. If your machine doesn't find them, install KoPubWorld
  % Dotum (Medium / Bold / Light) into ~/Library/Fonts/ and rebuild the
  % font cache, or comment this block out and uncomment the Apple SD Gothic
  % Neo fallback below.
  \setCJKmainfont[%
    BoldFont=KoPubWorldDotumBold,
    AutoFakeSlant=0.18
  ]{KoPubWorldDotumMedium}
  \setCJKsansfont[%
    BoldFont=KoPubWorldDotumBold,
    AutoFakeSlant=0.18
  ]{KoPubWorldDotumMedium}
  \setCJKmonofont{KoPubWorldDotumLight}

  % --- Apple SD Gothic Neo fallback (uncomment if KoPub unavailable) --------
  % \setCJKmainfont{Apple SD Gothic Neo}
  % \setCJKsansfont{Apple SD Gothic Neo}
  % \setCJKmonofont{Apple SD Gothic Neo}

  % --- Pretendard fallback (uncomment if installed) -------------------------
  % \setCJKmainfont{Pretendard}
  % \setCJKsansfont{Pretendard}
  % \setCJKmonofont{Pretendard}

  % Slightly looser line spread improves Hangul readability on slides
  \linespread{1.15}
\fi

% -----------------------------------------------------------------------------
% PALETTE — Beamer theme + TikZ-snippet colors.
% Load xcolor and define colors BEFORE anything that references them
% (notably hyperref, hypersetup, and the Beamer color assignments).
% -----------------------------------------------------------------------------
\usepackage{xcolor}

\definecolor{primary-blue}{HTML}{012169}      % headings, accents
\definecolor{primary-gold}{HTML}{B9975B}      % emphasis, borders
\definecolor{highlight-yellow}{HTML}{F2A900}  % markers, alerts
\definecolor{light-bg}{HTML}{E8EDF5}          % title / section backgrounds
\definecolor{jet}{HTML}{1A1A1A}               % body text

% Semantic colors (used in diagrams and annotations)
\definecolor{positive}{HTML}{15803D}          % good / observed / identified
\definecolor{negative}{HTML}{B91C1C}          % bad / problematic / bias
\definecolor{neutral}{HTML}{525252}           % reference / context

% Highlight accents (match .hi-* classes in the SCSS)
\definecolor{hi-slate}{HTML}{314F4F}
\definecolor{hi-green}{HTML}{2E8B57}
\definecolor{hi-red}{HTML}{C41E3A}

% Semantic aliases so skill templates and snippets can write
% \textcolor{accent}{...} without hard-coding "primary-blue".
\colorlet{accent}{primary-blue}
\colorlet{accent2}{primary-gold}

% -----------------------------------------------------------------------------
% Hyperref — loaded AFTER xcolor + palette so link colors resolve.
% -----------------------------------------------------------------------------
\usepackage{hyperref}
\hypersetup{colorlinks=true, linkcolor=primary-blue, urlcolor=primary-blue,
            citecolor=primary-blue, pdfborder={0 0 0}}

% -----------------------------------------------------------------------------
% Beamer theme assignments (only apply under Beamer).
%
% We detect Beamer via \@ifclassloaded{beamer}, which is the documented,
% non-internal way. This must be wrapped in \makeatletter/\makeatother
% because \@ifclassloaded contains an @-catcode character.
% -----------------------------------------------------------------------------
\makeatletter
\@ifclassloaded{beamer}{%
  \usefonttheme{professionalfonts}
  \setbeamercolor{structure}{fg=primary-blue}
  \setbeamercolor{frametitle}{fg=primary-blue}
  \setbeamercolor{title}{fg=primary-blue}
  \setbeamercolor{subtitle}{fg=primary-gold}
  \setbeamercolor{author}{fg=primary-blue}
  \setbeamercolor{institute}{fg=neutral}
  \setbeamercolor{date}{fg=primary-gold}
  \setbeamercolor{itemize item}{fg=primary-blue}
  \setbeamercolor{itemize subitem}{fg=primary-gold}
  \setbeamercolor{alerted text}{fg=primary-gold}
  \setbeamercolor{block title}{fg=white, bg=primary-blue}
  \setbeamercolor{block body}{bg=light-bg}

  % Minimal footer: page X / N only, no navigation clutter
  \setbeamertemplate{navigation symbols}{}
  \setbeamertemplate{footline}{%
    \hfill{\color{neutral}\footnotesize\insertframenumber/\inserttotalframenumber}\hspace{0.7em}\vspace{0.3em}}
}{}
\makeatother

% -----------------------------------------------------------------------------
% TikZ setup — libraries the snippet gallery uses
% -----------------------------------------------------------------------------
\usepackage{tikz}
\usetikzlibrary{arrows.meta, positioning, calc, decorations.pathreplacing,
                fit, shapes.geometric, backgrounds}

% Shared TikZ styles — snippets and hand-written diagrams can pick these up
% via `\begin{tikzpicture}[every node/.style={...}]` or by naming specific
% styles (e.g., `\node[dag-node] (X) at (0,0) {X};`).
\tikzset{
  % Node styles for causal DAGs and similar
  dag-node/.style = {
    draw=primary-blue, fill=light-bg, rounded corners=2pt,
    minimum width=1.2cm, minimum height=0.9cm, align=center, font=\small
  },
  decision-node/.style = {
    draw=primary-gold, fill=white, diamond, aspect=1.8,
    minimum width=1.8cm, minimum height=1.2cm, align=center, font=\small,
    inner sep=1pt
  },
  % Edge styles — observed vs counterfactual
  observed-edge/.style = {-{Latex[length=2.5mm]}, thick, draw=jet},
  counterfactual-edge/.style = {-{Latex[length=2.5mm]}, thick, draw=neutral, dashed},
  confound-edge/.style = {-{Latex[length=2.5mm]}, thick, draw=negative, dashed},
  % Data-point styles for plots
  observed-dot/.style = {fill=primary-blue, draw=primary-blue, circle, inner sep=1.2pt},
  counterfactual-dot/.style = {fill=white, draw=neutral, circle, inner sep=1.2pt}
}

% -----------------------------------------------------------------------------
% Convenience macros
% -----------------------------------------------------------------------------
\newcommand{\muted}[1]{\textcolor{neutral}{#1}}
\newcommand{\key}[1]{\textcolor{primary-gold}{\textbf{#1}}}
\newcommand{\good}[1]{\textcolor{positive}{#1}}
\newcommand{\bad}[1]{\textcolor{negative}{#1}}

% Standout frame macro: full-bleed dark background for section transitions.
% Expand: \transitionslide{Your transition title}
\newcommand{\transitionslide}[1]{%
  \begingroup
  \setbeamercolor{background canvas}{bg=primary-blue}
  \begin{frame}[plain]
    \centering
    \vfill
    {\usebeamerfont{title}\color{white}\Large #1\par}
    \vfill
  \end{frame}
  \endgroup
}

% with TEXINPUTS=../Preambles:$TEXINPUTS (the /compile-latex skill does this).
%
% Required font: KoPubWorld Dotum (Medium / Bold). Install from
% https://www.kopus.or.kr/biz/electronic/font.do (free for academic use).
% Fallback options documented under "Korean font setup" below.
% =============================================================================

% -----------------------------------------------------------------------------
% Engine detection + encoding
%
% XeLaTeX is required (Korean rendering via xeCJK depends on it). Under pdfTeX
% xeCJK is unavailable, so we fall back to inputenc and emit a clear warning.
% -----------------------------------------------------------------------------
\usepackage{iftex}
\ifPDFTeX
  \usepackage[utf8]{inputenc}
  \PackageWarningNoLine{header}{%
    pdfTeX detected --- Korean text will not render. Compile with XeLaTeX.}
\fi

\usepackage{amsmath, amssymb, amsthm}
\usepackage{booktabs}
\usepackage{graphicx}

% -----------------------------------------------------------------------------
% Korean font setup (XeLaTeX only)
%
% Loads xeCJK and KoPubWorld Dotum. The fonts live under ~/Library/Fonts/ on
% macOS. We pass file names (not family names) because the family name is
% registered as Korean ("KoPubWorld돋움체") and fontspec is more robust with
% explicit file references.
%
% Fallback ladder if KoPubWorld Dotum is unavailable: comment out the block
% below and uncomment ONE of the alternatives.
% -----------------------------------------------------------------------------
\ifXeTeX
  \usepackage{xeCJK}
  \xeCJKsetup{%
    CJKspace=true,           % automatic CJK-Latin spacing
    AutoFakeSlant=0.18,      % italic-style fallback when no italic cut exists
    CJKecglue={\hskip 0.18em plus 0.04em minus 0.04em}%
  }

  % --- KoPubWorld Dotum (primary) -------------------------------------------
  % macOS Font Book registers these by PostScript name (no spaces). XeLaTeX's
  % fontspec finds them via the system font registry without needing a
  % Path= override. If your machine doesn't find them, install KoPubWorld
  % Dotum (Medium / Bold / Light) into ~/Library/Fonts/ and rebuild the
  % font cache, or comment this block out and uncomment the Apple SD Gothic
  % Neo fallback below.
  \setCJKmainfont[%
    BoldFont=KoPubWorldDotumBold,
    AutoFakeSlant=0.18
  ]{KoPubWorldDotumMedium}
  \setCJKsansfont[%
    BoldFont=KoPubWorldDotumBold,
    AutoFakeSlant=0.18
  ]{KoPubWorldDotumMedium}
  \setCJKmonofont{KoPubWorldDotumLight}

  % --- Apple SD Gothic Neo fallback (uncomment if KoPub unavailable) --------
  % \setCJKmainfont{Apple SD Gothic Neo}
  % \setCJKsansfont{Apple SD Gothic Neo}
  % \setCJKmonofont{Apple SD Gothic Neo}

  % --- Pretendard fallback (uncomment if installed) -------------------------
  % \setCJKmainfont{Pretendard}
  % \setCJKsansfont{Pretendard}
  % \setCJKmonofont{Pretendard}

  % Slightly looser line spread improves Hangul readability on slides
  \linespread{1.15}
\fi

% -----------------------------------------------------------------------------
% PALETTE — Beamer theme + TikZ-snippet colors.
% Load xcolor and define colors BEFORE anything that references them
% (notably hyperref, hypersetup, and the Beamer color assignments).
% -----------------------------------------------------------------------------
\usepackage{xcolor}

\definecolor{primary-blue}{HTML}{012169}      % headings, accents
\definecolor{primary-gold}{HTML}{B9975B}      % emphasis, borders
\definecolor{highlight-yellow}{HTML}{F2A900}  % markers, alerts
\definecolor{light-bg}{HTML}{E8EDF5}          % title / section backgrounds
\definecolor{jet}{HTML}{1A1A1A}               % body text

% Semantic colors (used in diagrams and annotations)
\definecolor{positive}{HTML}{15803D}          % good / observed / identified
\definecolor{negative}{HTML}{B91C1C}          % bad / problematic / bias
\definecolor{neutral}{HTML}{525252}           % reference / context

% Highlight accents (match .hi-* classes in the SCSS)
\definecolor{hi-slate}{HTML}{314F4F}
\definecolor{hi-green}{HTML}{2E8B57}
\definecolor{hi-red}{HTML}{C41E3A}

% Semantic aliases so skill templates and snippets can write
% \textcolor{accent}{...} without hard-coding "primary-blue".
\colorlet{accent}{primary-blue}
\colorlet{accent2}{primary-gold}

% -----------------------------------------------------------------------------
% Hyperref — loaded AFTER xcolor + palette so link colors resolve.
% -----------------------------------------------------------------------------
\usepackage{hyperref}
\hypersetup{colorlinks=true, linkcolor=primary-blue, urlcolor=primary-blue,
            citecolor=primary-blue, pdfborder={0 0 0}}

% -----------------------------------------------------------------------------
% Beamer theme assignments (only apply under Beamer).
%
% We detect Beamer via \@ifclassloaded{beamer}, which is the documented,
% non-internal way. This must be wrapped in \makeatletter/\makeatother
% because \@ifclassloaded contains an @-catcode character.
% -----------------------------------------------------------------------------
\makeatletter
\@ifclassloaded{beamer}{%
  \usefonttheme{professionalfonts}
  \setbeamercolor{structure}{fg=primary-blue}
  \setbeamercolor{frametitle}{fg=primary-blue}
  \setbeamercolor{title}{fg=primary-blue}
  \setbeamercolor{subtitle}{fg=primary-gold}
  \setbeamercolor{author}{fg=primary-blue}
  \setbeamercolor{institute}{fg=neutral}
  \setbeamercolor{date}{fg=primary-gold}
  \setbeamercolor{itemize item}{fg=primary-blue}
  \setbeamercolor{itemize subitem}{fg=primary-gold}
  \setbeamercolor{alerted text}{fg=primary-gold}
  \setbeamercolor{block title}{fg=white, bg=primary-blue}
  \setbeamercolor{block body}{bg=light-bg}

  % Minimal footer: page X / N only, no navigation clutter
  \setbeamertemplate{navigation symbols}{}
  \setbeamertemplate{footline}{%
    \hfill{\color{neutral}\footnotesize\insertframenumber/\inserttotalframenumber}\hspace{0.7em}\vspace{0.3em}}
}{}
\makeatother

% -----------------------------------------------------------------------------
% TikZ setup — libraries the snippet gallery uses
% -----------------------------------------------------------------------------
\usepackage{tikz}
\usetikzlibrary{arrows.meta, positioning, calc, decorations.pathreplacing,
                fit, shapes.geometric, backgrounds}

% Shared TikZ styles — snippets and hand-written diagrams can pick these up
% via `\begin{tikzpicture}[every node/.style={...}]` or by naming specific
% styles (e.g., `\node[dag-node] (X) at (0,0) {X};`).
\tikzset{
  % Node styles for causal DAGs and similar
  dag-node/.style = {
    draw=primary-blue, fill=light-bg, rounded corners=2pt,
    minimum width=1.2cm, minimum height=0.9cm, align=center, font=\small
  },
  decision-node/.style = {
    draw=primary-gold, fill=white, diamond, aspect=1.8,
    minimum width=1.8cm, minimum height=1.2cm, align=center, font=\small,
    inner sep=1pt
  },
  % Edge styles — observed vs counterfactual
  observed-edge/.style = {-{Latex[length=2.5mm]}, thick, draw=jet},
  counterfactual-edge/.style = {-{Latex[length=2.5mm]}, thick, draw=neutral, dashed},
  confound-edge/.style = {-{Latex[length=2.5mm]}, thick, draw=negative, dashed},
  % Data-point styles for plots
  observed-dot/.style = {fill=primary-blue, draw=primary-blue, circle, inner sep=1.2pt},
  counterfactual-dot/.style = {fill=white, draw=neutral, circle, inner sep=1.2pt}
}

% -----------------------------------------------------------------------------
% Convenience macros
% -----------------------------------------------------------------------------
\newcommand{\muted}[1]{\textcolor{neutral}{#1}}
\newcommand{\key}[1]{\textcolor{primary-gold}{\textbf{#1}}}
\newcommand{\good}[1]{\textcolor{positive}{#1}}
\newcommand{\bad}[1]{\textcolor{negative}{#1}}

% Standout frame macro: full-bleed dark background for section transitions.
% Expand: \transitionslide{Your transition title}
\newcommand{\transitionslide}[1]{%
  \begingroup
  \setbeamercolor{background canvas}{bg=primary-blue}
  \begin{frame}[plain]
    \centering
    \vfill
    {\usebeamerfont{title}\color{white}\Large #1\par}
    \vfill
  \end{frame}
  \endgroup
}

% with TEXINPUTS=../Preambles:$TEXINPUTS (the /compile-latex skill does this).
%
% Required font: KoPubWorld Dotum (Medium / Bold). Install from
% https://www.kopus.or.kr/biz/electronic/font.do (free for academic use).
% Fallback options documented under "Korean font setup" below.
% =============================================================================

% -----------------------------------------------------------------------------
% Engine detection + encoding
%
% XeLaTeX is required (Korean rendering via xeCJK depends on it). Under pdfTeX
% xeCJK is unavailable, so we fall back to inputenc and emit a clear warning.
% -----------------------------------------------------------------------------
\usepackage{iftex}
\ifPDFTeX
  \usepackage[utf8]{inputenc}
  \PackageWarningNoLine{header}{%
    pdfTeX detected --- Korean text will not render. Compile with XeLaTeX.}
\fi

\usepackage{amsmath, amssymb, amsthm}
\usepackage{booktabs}
\usepackage{graphicx}

% -----------------------------------------------------------------------------
% Korean font setup (XeLaTeX only)
%
% Loads xeCJK and KoPubWorld Dotum. The fonts live under ~/Library/Fonts/ on
% macOS. We pass file names (not family names) because the family name is
% registered as Korean ("KoPubWorld돋움체") and fontspec is more robust with
% explicit file references.
%
% Fallback ladder if KoPubWorld Dotum is unavailable: comment out the block
% below and uncomment ONE of the alternatives.
% -----------------------------------------------------------------------------
\ifXeTeX
  \usepackage{xeCJK}
  \xeCJKsetup{%
    CJKspace=true,           % automatic CJK-Latin spacing
    AutoFakeSlant=0.18,      % italic-style fallback when no italic cut exists
    CJKecglue={\hskip 0.18em plus 0.04em minus 0.04em}%
  }

  % --- KoPubWorld Dotum (primary) -------------------------------------------
  % macOS Font Book registers these by PostScript name (no spaces). XeLaTeX's
  % fontspec finds them via the system font registry without needing a
  % Path= override. If your machine doesn't find them, install KoPubWorld
  % Dotum (Medium / Bold / Light) into ~/Library/Fonts/ and rebuild the
  % font cache, or comment this block out and uncomment the Apple SD Gothic
  % Neo fallback below.
  \setCJKmainfont[%
    BoldFont=KoPubWorldDotumBold,
    AutoFakeSlant=0.18
  ]{KoPubWorldDotumMedium}
  \setCJKsansfont[%
    BoldFont=KoPubWorldDotumBold,
    AutoFakeSlant=0.18
  ]{KoPubWorldDotumMedium}
  \setCJKmonofont{KoPubWorldDotumLight}

  % --- Apple SD Gothic Neo fallback (uncomment if KoPub unavailable) --------
  % \setCJKmainfont{Apple SD Gothic Neo}
  % \setCJKsansfont{Apple SD Gothic Neo}
  % \setCJKmonofont{Apple SD Gothic Neo}

  % --- Pretendard fallback (uncomment if installed) -------------------------
  % \setCJKmainfont{Pretendard}
  % \setCJKsansfont{Pretendard}
  % \setCJKmonofont{Pretendard}

  % Slightly looser line spread improves Hangul readability on slides
  \linespread{1.15}
\fi

% -----------------------------------------------------------------------------
% PALETTE — Beamer theme + TikZ-snippet colors.
% Load xcolor and define colors BEFORE anything that references them
% (notably hyperref, hypersetup, and the Beamer color assignments).
% -----------------------------------------------------------------------------
\usepackage{xcolor}

\definecolor{primary-blue}{HTML}{012169}      % headings, accents
\definecolor{primary-gold}{HTML}{B9975B}      % emphasis, borders
\definecolor{highlight-yellow}{HTML}{F2A900}  % markers, alerts
\definecolor{light-bg}{HTML}{E8EDF5}          % title / section backgrounds
\definecolor{jet}{HTML}{1A1A1A}               % body text

% Semantic colors (used in diagrams and annotations)
\definecolor{positive}{HTML}{15803D}          % good / observed / identified
\definecolor{negative}{HTML}{B91C1C}          % bad / problematic / bias
\definecolor{neutral}{HTML}{525252}           % reference / context

% Highlight accents (match .hi-* classes in the SCSS)
\definecolor{hi-slate}{HTML}{314F4F}
\definecolor{hi-green}{HTML}{2E8B57}
\definecolor{hi-red}{HTML}{C41E3A}

% Semantic aliases so skill templates and snippets can write
% \textcolor{accent}{...} without hard-coding "primary-blue".
\colorlet{accent}{primary-blue}
\colorlet{accent2}{primary-gold}

% -----------------------------------------------------------------------------
% Hyperref — loaded AFTER xcolor + palette so link colors resolve.
% -----------------------------------------------------------------------------
\usepackage{hyperref}
\hypersetup{colorlinks=true, linkcolor=primary-blue, urlcolor=primary-blue,
            citecolor=primary-blue, pdfborder={0 0 0}}

% -----------------------------------------------------------------------------
% Beamer theme assignments (only apply under Beamer).
%
% We detect Beamer via \@ifclassloaded{beamer}, which is the documented,
% non-internal way. This must be wrapped in \makeatletter/\makeatother
% because \@ifclassloaded contains an @-catcode character.
% -----------------------------------------------------------------------------
\makeatletter
\@ifclassloaded{beamer}{%
  \usefonttheme{professionalfonts}
  \setbeamercolor{structure}{fg=primary-blue}
  \setbeamercolor{frametitle}{fg=primary-blue}
  \setbeamercolor{title}{fg=primary-blue}
  \setbeamercolor{subtitle}{fg=primary-gold}
  \setbeamercolor{author}{fg=primary-blue}
  \setbeamercolor{institute}{fg=neutral}
  \setbeamercolor{date}{fg=primary-gold}
  \setbeamercolor{itemize item}{fg=primary-blue}
  \setbeamercolor{itemize subitem}{fg=primary-gold}
  \setbeamercolor{alerted text}{fg=primary-gold}
  \setbeamercolor{block title}{fg=white, bg=primary-blue}
  \setbeamercolor{block body}{bg=light-bg}

  % Minimal footer: page X / N only, no navigation clutter
  \setbeamertemplate{navigation symbols}{}
  \setbeamertemplate{footline}{%
    \hfill{\color{neutral}\footnotesize\insertframenumber/\inserttotalframenumber}\hspace{0.7em}\vspace{0.3em}}
}{}
\makeatother

% -----------------------------------------------------------------------------
% TikZ setup — libraries the snippet gallery uses
% -----------------------------------------------------------------------------
\usepackage{tikz}
\usetikzlibrary{arrows.meta, positioning, calc, decorations.pathreplacing,
                fit, shapes.geometric, backgrounds}

% Shared TikZ styles — snippets and hand-written diagrams can pick these up
% via `\begin{tikzpicture}[every node/.style={...}]` or by naming specific
% styles (e.g., `\node[dag-node] (X) at (0,0) {X};`).
\tikzset{
  % Node styles for causal DAGs and similar
  dag-node/.style = {
    draw=primary-blue, fill=light-bg, rounded corners=2pt,
    minimum width=1.2cm, minimum height=0.9cm, align=center, font=\small
  },
  decision-node/.style = {
    draw=primary-gold, fill=white, diamond, aspect=1.8,
    minimum width=1.8cm, minimum height=1.2cm, align=center, font=\small,
    inner sep=1pt
  },
  % Edge styles — observed vs counterfactual
  observed-edge/.style = {-{Latex[length=2.5mm]}, thick, draw=jet},
  counterfactual-edge/.style = {-{Latex[length=2.5mm]}, thick, draw=neutral, dashed},
  confound-edge/.style = {-{Latex[length=2.5mm]}, thick, draw=negative, dashed},
  % Data-point styles for plots
  observed-dot/.style = {fill=primary-blue, draw=primary-blue, circle, inner sep=1.2pt},
  counterfactual-dot/.style = {fill=white, draw=neutral, circle, inner sep=1.2pt}
}

% -----------------------------------------------------------------------------
% Convenience macros
% -----------------------------------------------------------------------------
\newcommand{\muted}[1]{\textcolor{neutral}{#1}}
\newcommand{\key}[1]{\textcolor{primary-gold}{\textbf{#1}}}
\newcommand{\good}[1]{\textcolor{positive}{#1}}
\newcommand{\bad}[1]{\textcolor{negative}{#1}}

% Standout frame macro: full-bleed dark background for section transitions.
% Expand: \transitionslide{Your transition title}
\newcommand{\transitionslide}[1]{%
  \begingroup
  \setbeamercolor{background canvas}{bg=primary-blue}
  \begin{frame}[plain]
    \centering
    \vfill
    {\usebeamerfont{title}\color{white}\Large #1\par}
    \vfill
  \end{frame}
  \endgroup
}
